% =============================================================================
% Template: Trabalho Acadêmico (NBR 14724:2024)
%
% Compilação local (a partir da raiz do repositório):
%   TEXINPUTS=./src: latexmk -pdf templates/trabalho-academico/main.tex
% =============================================================================

\documentclass{abnt-trabalho-academico}

\usepackage{abnt-resumo}
\usepackage{abnt-citacoes}   % carrega abnt-refs automaticamente

\addbibresource{refs.bib}

% --- Metadados (preencha com seus dados) ---
\titulo{Título do Trabalho}
\subtitulo{subtítulo, se houver}
\autor{Nome Completo do Autor}
\orientador{Prof. Dr. Nome do Orientador}
% \coorientador{Prof. Dr. Nome do Coorientador}
\instituicao{Nome da Instituição}
\local{Cidade}
\ano{2026}
\natureza{Dissertação apresentada ao Programa de Pós-Graduação
  em X da Universidade Y como requisito parcial para obtenção
  do título de Mestre em X.}
\programa{Programa de Pós-Graduação em X}
% \area{Área de concentração}

% =============================================================================
\begin{document}

% --- Pré-textuais ---
\imprimircapa
\imprimirfolhaderosto

\dataaprovacao{1 de janeiro de 2026}
\membrodabanca{Prof. Dr. Membro Um}{Instituição Um}
\membrodabanca{Prof. Dr. Membro Dois}{Instituição Dois}
\membrodabanca{Prof. Dr. Membro Três}{Instituição Três}
\imprimirfolhadeaprovacao

\imprimirdedicatoria{%
  Dedicatória (opcional).}

\imprimiragradecimentos{%
  Agradecimentos (opcional).}

\imprimirepigrafe{%
  ``Citação inspiradora'' (opcional).}

\imprimirresumo{%
  Resumo em português. Deve conter entre 150 e 500 palavras para trabalhos
  acadêmicos, conforme NBR 6028:2021.%
}{Palavra-chave 1. Palavra-chave 2. Palavra-chave 3}

\imprimirabstract{%
  Abstract in English. Should contain between 150 and 500 words.%
}{Keyword 1. Keyword 2. Keyword 3}

\imprimirsumario

% --- Textuais ---
\chapter{Introdução}

Este é um template para trabalhos acadêmicos conforme a
ABNT NBR 14724:2024 \cite{eco1977}.

\chapter{Desenvolvimento}

Texto do desenvolvimento \cite{medeiros2019}.

\chapter{Conclusão}

Texto da conclusão.

% --- Pós-textuais ---
\postextual
\printbibliography

\end{document}
